\ficha{Flores Zendejas (2016), Financial markets, international organizations and conditional lending}{floreszendejas2016}

\begin{identificacao}
FLORES ZENDEJAS, Juan H. Financial markets, international organizations and
conditional lending: a long-term perspective. In: MALLARD, Grégoire; SGARD,
Jérôme (org.). \emph{Contractual knowledge}: one hundred years of legal
experimentation in global markets. Cambridge: Cambridge University Press, 2016.
p. 61-91. (Cambridge Studies in Law and Society). DOI 10.1017/CBO9781316442876.002.

Capítulo de livro, 31 páginas, lido integralmente, em inglês. Versão lida: a
cópia de acesso aberto depositada no Archive ouverte da Universidade de Genebra,
idêntica à publicada, com uma capa de repositório antes da p. 61.
\end{identificacao}

\bloco{Problema e tese}
O capítulo pergunta de onde vem o empréstimo condicionado e quanto as
organizações financeiras internacionais de hoje inovaram frente ao século XIX. A
resposta é que a condicionalidade nasceu no mercado primário de Londres e
repousava sobre dois atores: uma autoridade institucional que fixava os
requisitos de acesso à liquidez e intermediários financeiros com capacidade de
excluir de fato o devedor que não cumprisse o contrato (p. 62). A Liga das
Nações e o FMI legitimaram e difundiram essas práticas sem alterar os princípios
montados no século XIX. Daí decorre a tese avaliativa do texto: as organizações
internacionais devem ser julgadas pelo papel ex ante de facilitar o
desenvolvimento do mercado de dívida, não apenas pela gestão ex post das crises
(p. 63).

\bloco{Argumento}
\begin{enumerate}
\item A arquitetura oitocentista tinha um eixo institucional e um eixo de
mercado. A Bolsa de Londres podia vetar a listagem de títulos de governo em
default que não tivesse chegado a acordo satisfatório com os portadores, e assim
funcionava como corte arbitral e antecedente das cláusulas de ação coletiva
(p. 64). Essa cláusula de não default era a condição institucional genérica,
válida para todos.

\item Condição explícita sobre política econômica em empréstimo novo era rara. O
enforcement era induzido pelo mercado e dependia de concentração: os Rothschild
detinham em média cerca de 60\% do mercado no início do século e por isso
puderam impor, no empréstimo belga de 1832, uma cláusula política (p. 65).
Depois disso, o autor afirma que não houve outra tentativa séria até 1898
(p. 66). No conjunto, três de pelo menos noventa e três bancos emissores
tentaram condicionar seus empréstimos, e esses três cobriam 66\% do mercado
entre 1877 e 1895 (p. 70). Os emissores menores não assumiam o papel de gestores
de default, o que empurrou os investidores para os comitês de portadores, entre
eles a Corporation of Foreign Bondholders (p. 70).

\item Os três resgates condicionais da década de 1890 são o núcleo comparativo:
o comitê Rothschild para a Argentina em 1890, o empréstimo Hambro para a Grécia
em 1893 e o funding loan brasileiro de 1898, organizado pelos Rothschild
(p. 68). Todos visavam à solvência do devedor e todos, diz o autor, deram
resultados decepcionantes, porque a posição financeira desses governos nunca se
estabilizou. O funding loan é definido em nota como pagamento dos juros da
dívida externa com novos títulos equivalentes ao devido por prazo limitado, em
geral três anos (p. 68).

\item Os três casos diferem no grau de intrusão. A Argentina teve moratória de
três anos e desconto renegociado em 1893, precedidos de relatório de uma
comissão de banqueiros apresentado ao Banco da Inglaterra. O General Bond
enumerava compromissos que eram as recomendações dos banqueiros: vincular o
empréstimo à receita aduaneira cobrada em pesos ouro, reduzir o meio circulante
e formar uma Caja de Conversión que levaria à adoção do padrão-ouro, e não
contrair nova dívida externa (p. 68). A Grécia teve condições sobre receitas
dadas em garantia e uma Caisse de agentes bancários encarregada do empréstimo,
cláusulas que a convenção final excluiu e transferiu para decreto do Executivo.
O acordo não foi ratificado e o país entrou em default. A solução veio em 1896
com governos credores e uma Comissão Internacional de Controle, criada em 1897
após a derrota na guerra greco-turca, que também recomendou equilibrar o
orçamento e revalorizar o dracma por contração monetária (p. 69, 70, 71). O
Brasil de 1898 foi, na formulação do autor, sucesso relativo de cumprimento, com
uma única condição de política econômica, contrair o meio circulante para
valorizar a moeda (p. 69).

\item A tipologia das condições explica a diferença de eficácia. As condições
sobre gestão dos fundos e fiscalização das receitas dadas em garantia eram
contratuais. As que buscavam reforçar a capacidade de pagamento dependiam só do
mecanismo de mercado. Babb e Carruthers classificam esse arranjo como a forma
mais estreita de condicionalidade, aquela em que credores privados são os únicos
a exigir condições (p. 69). Sem força legal, a opção de saída permanecia aberta.
A escada de enforcement disponível era o comitê de portadores, o veto da Bolsa
de Londres e, no extremo, a intervenção de governo credor, sendo a diplomacia
das canhoneiras evento raro e não a norma do mercado (p. 70).

\item A Liga das Nações reproduziu o procedimento dos banqueiros privados e
acrescentou o que faltava a eles: apoio de governos credores e um enforcement
mais amplo, com protocolos, agentes fiscais permanentes e um Comissário-Geral. O
conteúdo das políticas promovidas era comércio liberal, orçamento equilibrado,
adoção do padrão-ouro e bancos centrais autônomos (p. 79). No caso austríaco de
1923, Baring e Rothschild só aceitaram subscrever depois que a Liga montou o
plano de longo prazo (p. 73). A fragilidade estava no descasamento de prazos:
cumprido o programa, o incentivo de mercado desaparecia muito antes do
vencimento do empréstimo, e nos anos 1930 a maioria dos League loans entrou em
default.

\item O FMI herdou o efeito de certificação, mas dispõe de recursos próprios,
age de modo anticíclico e pode operar como emprestador de última instância, o
que o desloca da condicionalidade no mercado primário para a gestão de crises
depois de 1982.
\end{enumerate}

\chave{no other serious attempt to explicitly require conditions on new loans were included until 1898}{p. 66}

\chave{the sole condition regarding economic policy was reducing the money supply in an attempt to allow the local currency to appreciate}{p. 69}

\bloco{Conceitos e definições operacionais}
Condicionalidade: implementação de políticas específicas para obter assistência
financeira, definição que o autor toma de Dreher (p. 61). Cláusula de não
default: condição institucional genérica imposta pela Bolsa de Londres, que
suspende o acesso ao mercado primário até haver acordo com os portadores
(p. 64). Enforcement induzido pelo mercado: capacidade do banco emissor
dominante de fixar termos, montante e calendário de liberação dos fundos, sem
recurso judicial. Opção de saída: alternativa de financiamento disponível ao
devedor, que mede o custo real de descumprir. Money doctor: o agente externo que
prescreve e monitora o programa, papel exercido por bancos no século XIX e pela
Liga nos anos 1920. Efeito de certificação: o aval da organização como senha de
acesso ao crédito privado. O padrão-ouro aparece aqui não como objeto autônomo,
e sim como um dos itens do receituário liberal difundido pela Liga (p. 79) e
como resultado esperado do programa argentino de 1890 (p. 68).

\bloco{Evidência e método}
Método narrativo comparativo, sem modelo formal e sem séries quantitativas
próprias. As fontes são arquivos bancários e institucionais: Banco da
Inglaterra, ING Baring, Hambro, Arquivos Rothschild, Arquivo da Liga das Nações,
Arquivos do FMI e Morgan Library. O autor cita contratos, prospectos, resoluções
do Conselho da Liga e correspondência de banqueiros e diplomatas. Os números de
concentração de mercado vêm de Flandreau e coautores, não de reconstrução
própria. O recorte vai do início do século XIX à crise mexicana de 1982 e seus
desdobramentos.

\bloco{Agentes e vertentes}
Bancos: Rothschild, Baring, Hambro, J.P. Morgan. Instituições: Bolsa de Londres,
Corporation of Foreign Bondholders, Banco da Inglaterra, Foreign Office,
Organização Econômica e Financeira da Liga, Delegação do Ouro da Liga, FMI.
Historiograficamente, o autor se opõe aos que julgam a Liga pelo fracasso ex
post, como Pauly, Schuker e Fior, e se aproxima de Sargent e de Santaella, que
veem êxito relativo em queda de inflação e ganho de credibilidade. Contra
Mitchener e Weidenmier, alinha-se a Tomz e a seus próprios trabalhos ao tratar a
coerção militar como exceção. Contra Pauly, sustenta que o laissez-faire
britânico valia para o regime monetário, o padrão-ouro, mas não para o mercado
financeiro (p. 71).

\bloco{Críticas e limites}
O caso brasileiro ocupa poucas linhas e não repousa em pesquisa de arquivo
própria, remetendo a Flandreau e Flores. O sucesso relativo de cumprimento de
1898 é afirmado sem métrica de cumprimento e sem discussão do custo interno da
contração monetária. O capítulo não trata da Caixa de Conversão, não menciona o
Convênio de Taubaté e não diz nada sobre o debate público brasileiro. A menção
ao segundo resgate de 1914, também organizado pelos Rothschild, fica em uma
frase, sem análise das condições (p. 69). O autor exclui explicitamente do
escopo a relação entre dívida soberana e expansão colonial (p. 67), o que
enfraquece a tese de que a coerção era rara. A avaliação positiva da Liga é
normativa e o texto a assume como tal (p. 63). Nada disso se transfere sem
mediação para 1906: o capítulo oferece o enquadramento da condicionalidade, não
a interpretação do caso brasileiro.

\begin{dialogo}
\item Procurar, nos editoriais de 1906, a memória do funding loan de 1898 e de
sua única condição, a contração do meio circulante. Vocabulário: ``funding'',
``Rothschild'', ``compromissos com os credores'', ``resgate do papel-moeda'',
``incineração''. Testa se o enforcement que o autor descreve como induzido pelo
mercado havia sido internalizado pelos agentes brasileiros, ou se em 1906 já era
tratado como obrigação extinta.
\item Verificar se a defesa da estabilização e da taxa baixa aparece como
escolha nacional ou como exigência de credor. Se a imprensa ortodoxa invocar
Londres, isso corrobora a condicionalidade de mercado descrita no capítulo. Se
recorrer apenas a argumento técnico, sugere que o constrangimento externo já não
precisava ser nomeado.
\item Rastrear quem os jornais nomeiam como fiscal do compromisso: banqueiros,
portadores de títulos, ``praça de Londres'' ou governo inglês. O capítulo
sustenta que o enforcement era bancário e que a intervenção de governo credor
era exceção (p. 70). A distribuição das menções testa isso do lado brasileiro.
\item Buscar como circula a comparação argentina em 1906, se a Caja de
Conversión aparece como modelo técnico ou como peça de um programa de credores,
já que no General Bond de 1890 ela é compromisso do governo argentino perante os
banqueiros (p. 68). Alimenta a parte IV do capítulo.
\item Procurar invocação dos precedentes grego e turco, controle externo e
comissão internacional, como ameaça retórica no debate brasileiro. Testa se a
escada de enforcement descrita no capítulo era conhecida e usada como argumento
pelos publicistas.
\end{dialogo}

\usonocapitulo{Parte II. Serve para situar o funding loan de 1898 numa classe
pequena e datada de resgates condicionais, três em toda a década de 1890, e para
sustentar duas afirmações: que o enforcement dos Rothschild era de mercado e não
jurídico, e que a única condição de política econômica imposta ao Brasil foi
contrair o meio circulante para valorizar o mil-réis (p. 69). Também fornece o
contraste institucional que separa o caso brasileiro do grego, onde houve
comissão internacional de controle, e a ponte com a parte IV pela Caja de
Conversión argentina como compromisso contratual de 1890 (p. 68). Como o
capítulo não trata da Caixa de Conversão nem do debate de imprensa, ele entra
como moldura do constrangimento externo, não como evidência sobre 1906.}
